\documentclass[10pt,a4paper]{amsart}
\usepackage[margin=19mm]{geometry}
\usepackage{amsmath,amssymb,mathtools}
\usepackage[scr=rsfs,cal=euler]{mathalfa}
\usepackage{xfrac}
\usepackage[unicode,hidelinks,bookmarks=false]{hyperref}
\newcommand{\ie}{i.e.\ }
\newcommand{\cf}{cf.\ }
\newcommand{\resp}{respectively\ }
\newcommand{\Subsection}{\subsection}
\providecommand{\nobreakdash}{\nobreak\hbox{-}}
\setlength{\emergencystretch}{2em}
\multlinegap0pt
\usepackage{bm}
\usepackage{bbm}
\usepackage{dsfont}
\usepackage{xspace}
\usepackage{cancel}
\usepackage{mathtools}
\usepackage[shortlabels]{enumitem}
\usepackage{amsthm}
\makeatletter
\@ifundefined{theorem}{\newtheorem{theorem}{Theorem}}{}
\makeatother
\makeatletter
\expandafter\ifx\csname claim\endcsname\relax
\newenvironment{claim}[1]{\par\noindent\underline{Claim:}\space#1}{}
\fi
\newcommand{\hodge}{\!\star\!}
\newcommand{\iprod}{\mathbin{\lrcorner}}
\makeatother
\title[Spectral geometry of the curl operator on smoothly bounded d]{Spectral geometry of the curl operator on smoothly bounded domains}
\author{Josef Greilhuber, Willi Kepplinger}
\date{Source: arXiv:2502.13067v2 (CC BY 4.0)}
\begin{document}
\maketitle

\noindent\emph{Source and licence.} Spectral geometry of the curl operator on smoothly bounded domains. Version arXiv:2502.13067v2 (CC BY 4.0), distributed under the Creative Commons Attribution 4.0 International licence, as recorded in the arXiv licence metadata for this exact version and shown on the abstract page https://arxiv.org/abs/2502.13067v2. Full rights notice: licenses/arxiv-2502.13067-CC-BY-4.0.txt.\par\smallskip

\noindent\textit{Editorial scope.} Counted unit: the Theorem whose source label is \texttt{thm:hadamard\_rule} in the article (source0.tex lines 425--440, source label \texttt{thm:hadamard\_rule}), delivered together with the article's own complete original proof of it (source0.tex lines 441--518, one \texttt{proof} environment of 78 lines). No mathematics is written, repaired or re-derived: statement and proof are the article's own text line for line. Required context retained uncounted: none. Every definition, display and label of those lines is retained unchanged. Omitted from this package and not redistributed: 1--424, 519--738. Nothing inside a retained excerpt is omitted or altered: every retained body is delivered exactly as the article's lines are written, so no omission receipt of any kind accompanies this package. Citations: the retained excerpts carry no citation command at all, so no bibliography entry of the article is transcribed and no reference is redistributed. Reference adaptation: none was needed and none was made; every reference inside the delivered unit resolves inside it, so no unresolved reference or citation is produced. Printed slips kept literal: no printed word, symbol, hypothesis or number of the counted statement or of its proof was altered. Layout and class: the article's own class is \texttt{amsart} with packages including graphicx, amsmath, amssymb, amsthm, url, color, tikz, tikz-cd, dsfont, comment, cleveref, enumitem, thmtools, thm-restate; the wrapper is the lane's pinned \texttt{amsart} editorial wrapper at 10pt on A4, which restates the theorem-like environments the retained text needs and adds the article's own macro definitions that the retained text needs (\texttt{\textbackslash hodge}, \texttt{\textbackslash iprod}), with extra wrapper packages (bm, bbm, dsfont, xspace, cancel, mathtools, enumitem). The delivered numbering is the unit's own. Assets: no retained span contains a figure environment, a graphics-inclusion command, a PDF-page inclusion, a table or a TikZ picture; no image file is copied into this package, the package declares no asset dependency and no image file is redistributed with it. Licence: version arXiv:2502.13067v2 of this article is distributed under the Creative Commons Attribution 4.0 International licence, as recorded in the arXiv licence metadata for this exact version and shown on the abstract page https://arxiv.org/abs/2502.13067v2. A full rights notice is delivered at licenses/arxiv-2502.13067-CC-BY-4.0.txt. Characters: every retained excerpt is pure ASCII, so the delivered engine, XeLaTeX with its default Latin Modern text font, reports no missing character.
\begin{theorem}[Hadamard rule]
\label{thm:hadamard_rule}
Let $D \subseteq \mathbb R^3$ be a smoothly bounded domain and $\Phi_t: D \to \mathbb R^3$ a smooth one-parameter family of embeddings of $D$. Write $D_t = \Phi_t (D)$, and $X_t = \frac{d\Phi_t}{dt}$. Furthermore, fix a Lagrangian $L \subseteq H^1_{dR}(\partial D;\mathbb C)$.
\begin{enumerate}[label=(\alph*)]
    \item Suppose that $u_t$ is a smooth $1$-parameter family of closed, boundary-parallel $2$-forms which solves $d ( \star u_t ) = \lambda_t u_t$ on $D_t$ and satisfies $[\Phi_t^\ast ( \star u_t)] \in L$ for all $t$.
    Then
    \begin{align}
        \Dot{\lambda} = - \lambda \int_{\partial D} (X \cdot \nu) \, |u|^2\, d\sigma,
    \end{align}
    where $\nu$ is the outward pointing unit normal of $\partial D$ and $d\sigma$ is the induced surface measure on $\partial D$.
    \item Let $v_t$ be another smooth $1$-parameter family of closed, boundary-parallel $2$-forms solving $d ( \star v_t ) = \mu_t v_t$ satisfying $[\Phi_t^\ast ( \star v_t)] \in L$ for all $t$. Suppose $\mu_0 = \lambda_0 \neq 0$ and $\left( u_t,v_t\right)_{L^2(D_t)} = 0$ for all $t$. Then 
    \begin{align}
        0 = \int_{\partial D} (X \cdot \nu) \, (u \cdot v) \, d\sigma.
    \end{align}
\end{enumerate}
\end{theorem}

\begin{proof}
We will first prove the identities \eqref{identity 1}--\eqref{identity 3} below, which follow from $L^2$-normalization, Lagrangian boundary conditions and parallel boundary conditions, respectively. In these, $w_t$ is a $1$-parameter family of eigenforms which may stand for both $u_t$ and $v_t$. We write $\rho_t$ as a placeholder for $\lambda_t$ and $\mu_t$, so that $d ( \star w_t ) = \rho_t w_t$. Then
\begin{align}\label{identity 1}
    \int_{\partial D} (X\cdot \nu)\,( \star u \cdot \star w)\,d\sigma + \int_{D} \star \Dot{u}\wedge \bar w +\int_{D} u\wedge \star \Dot{\bar w} =0
\end{align}
\begin{align}\label{identity 2}
    \int_{\partial D}\star\dot{u}\wedge\star \bar w= - \lambda \int_{\partial D} (X\iprod u) \wedge \star \bar w
\end{align}
\begin{align}\label{identity 3}
    \int_{\partial D} (X \iprod u) \wedge \star \bar w=\int_{\partial D} (X\cdot \nu) (u \cdot w)\,d\sigma
\end{align}

Identity (\ref{identity 1}) follows from differentiating the expression
\begin{align*}
\langle u_t,w_t\rangle_{D_t} = \int_{ D_t}u_t\wedge \star \bar w_t = \int_{D_t} (\star u_t \cdot \star w_t) \ dV_{\mathbb R^3},
\end{align*}
which is constantly $1$ if $w_t = u_t$ or constantly $0$ if $w_t = v_t$. Thus,
\begin{align*}
    0 = \frac{d}{dt}\lvert_{t=0} \int_{ D_t} u_t \wedge \star \bar w_t &= \int_{\partial D} (X \cdot \nu) (\star u \cdot \star w) + \int_{D} \dot u \wedge \star \bar w + \int_D u \wedge \star \dot{\bar w} \\
    &= \int_{\partial D} (X \cdot \nu) (\star u \cdot \star w) + \int_{D} \star \dot u \wedge \bar w + \int_D u \wedge \star \dot{\bar w},
\end{align*}
noting that $\int_D \dot u \wedge \star {\bar w} = \int_D (\star \dot u, \star w) \, dV_{\mathbb R^3} = \int_D \star \dot u \wedge \bar w$.


We begin the proof of \eqref{identity 2} by differentiating 
\begin{align*}
    \frac{d}{dt}\lvert_{t=0}\Phi_t^\ast (\star u_t)= \star \dot{u} + X\iprod (d\hodge u) + d(X\iprod \star u)
\end{align*}
Since $[\star w]$ and $[\Phi^\ast_t (\star u_t)]$ both belong to the same Lagrangian $L \subseteq H^1_{dR}(\partial D; \mathbb C)$,
\begin{align*}
    0 = \int_{\partial D} \Phi_t^\ast (\star u_t) \wedge \star \bar w.
\end{align*}
Differentiating this identity at $t=0$ and applying Stokes' theorem on $\partial D$ yields
\begin{align*}
    0 &= \int_{\partial D} \star \dot u \wedge \star \bar w+ \lambda \int_{\partial D} (X \iprod u) \wedge \star \bar w  + \int_{\partial D} d ( X \iprod \star u) \wedge \star \bar w  \\
    &= \int_{\partial D} \star \dot u \wedge \star \bar w + \lambda \int_{\partial D} (X \iprod u) \wedge \star \bar w  - \int_{\partial D} (X \iprod \star u) \wedge (d \hodge \bar w)\\
    &= \int_{\partial D} \star \dot u \wedge \star \bar w+ \lambda \int_{\partial D} (X \iprod u) \wedge \star \bar w,
\end{align*}
where the last term in the second line vanished since $\iota_{\partial D}^\ast(d \hodge \bar w) = \rho \, \iota_{\partial D}^\ast \bar w = 0$. 





For identity (\ref{identity 3}), let  $X^\perp = (X\cdot \nu)\, \nu$ and $X^T = X-X^\perp$. The graded product rule for the interior product yields
\begin{align*}
    (X \iprod u) \wedge \star \bar w&= 
    X \iprod (u \wedge \star \bar w) - u \wedge (X \iprod \star \bar w) \\
    &= X^\perp \iprod (u \wedge \star \bar w) + X^T \iprod (u \wedge \star \bar w) - u \wedge (X \iprod \star \bar w)
\end{align*}
When pulled back to $\partial D$, the terms $X^T \iprod ( u \wedge \star \bar w)$ and $u \wedge (X \iprod \star \bar w)$ both vanish -- the former because $X^T$ is parallel to $\partial D$ and the latter because $\iota_{\partial D}^\ast u = 0$. Hence,
\begin{align*}
    \int_{\partial D} (X \iprod u) \wedge \star \bar w & = \int_{\partial D} X^\perp \iprod (u \wedge \star \bar w) \\
    &= \int_{\partial D} (X\cdot \nu)\, ( u \cdot w) \, (\nu \iprod dV_{\mathbb R^3}) \\
    & = \int_{\partial D} (X\cdot \nu)\, ( u \cdot w)\,d\sigma
\end{align*}

With these identities in hand we compute
\begin{align*}
    \Dot{\lambda}&= \frac{d}{dt}\lvert_{t=0} \int_{ D_t}d(\star u_t)\wedge \star \bar u_t \\
    &= \lambda \int_{\partial  D}\big(X\cdot \nu\big) \, \lvert \star u\rvert^2 \, d\sigma + \int_{ D} d(\star \Dot{u})\wedge \star \bar u + \lambda \int_{ D} u\wedge \star \Dot{\bar u} \\
    &=\lambda \int_{\partial D}\big(X \cdot \nu\big) \, \lvert \star u\rvert^2 \,d\sigma + \int_{\partial  D} \star \Dot{u}\wedge \star \bar u + \lambda \int_D \star \dot u \wedge \bar u + \lambda \int_D u\wedge \star \Dot{\bar u}\\
    &\overset{(\ref{identity 1})}{=} \int_{\partial  D} \star \Dot{u} \wedge \star \bar u \overset{(\ref{identity 2})}{=} - \lambda \int_{\partial  D} (X \iprod u) \wedge \star\bar u \\
    &\overset{(\ref{identity 3})}{=} - \lambda\int_{ \partial D} (X\cdot \nu) \lvert u\rvert^2\,d\sigma,
\end{align*}
yielding our Hadamard-type rule, part (a). The third line resulted from an application of Stokes' theorem to $\int_{\partial D} \star \dot u \wedge \star \bar u$, keeping in mind that $d \hodge \bar u = \lambda \bar u$ since $\lambda$ is real.

For part (b) we essentially repeat the same computations, keeping in mind that $\mu_0=\lambda_0$.
\begin{align*}
    0&= \frac{d}{dt}\lvert_{t=0} \int_{ D_t}d(\star u_t)\wedge \star \bar v_t \\
    &= 
    \lambda \int_{\partial  D} \big(X\cdot \nu\big) \, (\star u \cdot \star v) \, d\sigma + \int_{ D} d(\star \Dot{u})\wedge \star \bar v + \lambda \int_{ D} u\wedge \star \Dot{\bar v} \\
    &= \lambda \int_{\partial  D} \big(X\cdot \nu\big) \, (\star u \cdot \star v) \, d\sigma + \int_{\partial  D} \star \Dot{u}\wedge \star \bar v + \mu \int_D \star \dot u\wedge \bar v + \lambda \int_D u\wedge \star \Dot{\bar v} \\
    &\overset{(\ref{identity 1})}{=} \int_{\partial  D} \star \Dot{u} \wedge \star \bar v \overset{(\ref{identity 2})}{=} - \lambda \int_{\partial  D} (X \iprod u) \wedge \star \bar v \\
    &\overset{(\ref{identity 3})}{=} - \lambda \int_{ \partial D} (X\cdot \nu) \, ( u \cdot v)\,d\sigma
\end{align*}
Note that no assumptions on the connectedness of $\partial D$ were made.
\end{proof}

\end{document}
